\chapter{Ejercicios}
\begin{observacion}
    Para comprobar si un operador lineal y continuo $T$ es compacto en $\ell_2$ podemos fiarnos por la siguiente regla. \\ \\
    \noindent
    En el espacio $\ell_2$, un operador que se limita a multiplicar cada componente por un número ($T(x) = (\alpha_1 x_1, \alpha_2 x_2, \dots)$) se llama operador diagonal. Para estos operadores existe una regla intuitiva infalible:
    
    \textit{Un operador diagonal es compacto si y solo si sus coeficientes tienden a cero cuando $k$ va al infinito ($\lim_{k \to \infty} \alpha_k = 0$)}
    \noindent
    Es claro que igualmente habrá que demostrarlo formalmente, pero nos da una idea clave para que hacer.
\end{observacion}

\section{Folio 1}
\begin{ejercicio}
    Sea $H$ un espacio con producto interno y $S$ un sistema ortonormal. Demostrar que las siguientes afirmaciones son equivalentes:
    \begin{itemize}
        \item[(i)] $S$ es completo
        \item[(ii)] $\forall x \in H, \quad x \perp S \implies x = 0$
    \end{itemize}
    \begin{description}
    \item[$i)\Longrightarrow ii))$] Sea $x \in \mathcal{H} : (x, s) = 0 \quad \forall s \in S$. Entonces, o bien $x=0$, o bien 
    $$\frac{1}{\|x\|}(x, s) = \left(\frac{x}{\|x\|}, s\right) = 0 \implies \frac{x}{\|x\|} \perp S$$
    Por lo tanto, $S \cup \{ \frac{x}{\|x\|} \}$ sería aún un sistema ortonormal, lo cual contradice la maximalidad de $S$ (a menos que $x=0$).
    \item[$ii)\Longrightarrow i)$] Si $x \perp S$, entonces $x = 0$, esto implica que no existe ningún sistema ortonormal $S' \supsetneq S$. Por lo tanto, $S$ es completo.
\end{description}
\end{ejercicio}

\begin{ejercicio} % //TODO: hacer
    Sea $H$ un espacio de Hilbert y $S$ un sistema ortonormal. Demostrar que las siguientes afirmaciones son equivalentes:
    \begin{itemize}
        \item[(i)] $S$ es completo
        \item[(ii)] $H = \overline{[S]}$, es decir, $[S]$ es denso en $H$ $\big([S] = [\text{span } S]\big)$
    \end{itemize}
\end{ejercicio}

\begin{ejercicio}
    Sea $H$ un espacio de Hilbert y $\{e_1, \dots, e_n, \dots\}$ un sistema ortonormal numerable. \\
    Demostrar que $e_n \rightharpoonup 0$, pero $\{e_n\}$ no converge fuertemente a ningún elemento de $H$ (en particular, $e_n \not\longrightarrow 0$ fuertemente). \\

    \noindent
    Recordemos la definición de convergencia débil, para $\mathcal{H}$ un espacio normado aparte de de Hilbert. Se dice que $e_n$ converge débilmente a $0$, si
    \begin{equation*}
        F(e_n) \longrightarrow F(0) = 0 \quad \forall F \in H^*
    \end{equation*}
    Como $\mathcal{H}$ es un espacio de Hilbert, por el teorema de representación de Riesz, que dice lo siguiente: \\

    \textit{ Sea H un espacio de Hilbert, entonces
    \begin{equation*}
        \forall F\in H^*, \quad \exists! y_F \in H \text{ tal que } F(x) = (x, y_F) \quad \forall x \in H
    \end{equation*}}
    \noindent
    entonces sabemos que cada funcional $F \in H^*$ se puede escribir como $F(x) = (x, y)$ para algún $y \in H$. Por lo tanto, la condición de convergencia débil se puede reescribir como
    \begin{equation*}
        (e_n, y) \longrightarrow (0, y) = 0 \quad \forall y \in H
    \end{equation*}
    Aplicamos la Desigualdad de Bessel al sistema ortonormal $\{e_n\}$ y al vector $y$
    $$\sum_{n=1}^{\infty} |F(e_n)|^2 = \sum_{n=1}^{\infty} |( e_n, y )|^2 \le \|y\|^2 < \infty$$
    Como la serie de números reales $\sum |F(e_n)|^2$ converge (su suma es menor que infinito), el término general debe tender a cero de forma obligatoria cuando $n \to \infty$
    $$\lim_{n \to \infty} |F(e_n)|^2 = 0 \implies \lim_{n \to \infty} F(e_n) = 0$$
    Como hemos partido de un funcional $F \in H^*$ absolutamente genérico y hemos demostrado que $F(e_n) \to 0 = F(0)$, aplicamos la Definición 5.1 de tus apuntes y concluimos
    $$e_n \rightharpoonup 0 $$
    Ahora vamos a probar que $\{e_n\}$ no converge fuertemente a ningún elemento de $H$. Para ello, veámos cual es la distancia entre dos elemenos cualesquiera del sistema ortonormal, es decir, dados $m,n\in \N$ con $m \neq n$, calculamos la norma de la diferencia entre $e_n$ y $e_m$
    $$\|e_n - e_m\|^2 = ( e_n - e_m, e_n - e_m )= (e_n, e_n) - (e_n, e_m) - (e_m, e_n) + (e_m, e_m) = 1 - 0 - 0 + 1 = 2$$
    entonces 
    \begin{equation*}
        \|e_n - e_m\| = \sqrt{2} \quad \forall m,n \in \mathbb{N}, m \neq n
    \end{equation*}
    Si recordamos la definición de sucesión de cauchye, tenemos \\
    
    \textit{Una sucesión es de Cauchy si, a medida que avanzamos hacia el infinito, los términos de la sucesión se van pegando cada vez más entre sí. Es decir, para cualquier distancia pequeña $\epsilon > 0$, existe un momento a partir del cual cualesquiera dos elementos $e_n, e_m$ están a una distancia menor que $\epsilon$:$$\lim_{n, m \to \infty} \|e_n - e_m\| = 0$$}

    \noindent
    pero nosotros tenemos que si tomamos dos elementos distintos cualesquiera de nuestro sistema, la distancia entre ellos siempre es una constante fija de $\sqrt{2}$, da igual lo grandes que sean $n$ y $m$, por lo que
    $$\lim_{n, m \to \infty} \|e_n - e_m\| = \sqrt{2} \neq 0$$
    lo que significa que $\{e_n\}$ no es una sucesión de Cauchy, y por lo tanto no es una sucesión convergente en un espacio de Hilbert (que es completo). En particular, no puede converger a $0$ ni a ningún otro elemento de $H$.
\end{ejercicio}

\begin{ejercicio} % //TODO: hacer
    Sea $X = C^0([0,1])$ con la norma usual $\|\cdot\|_\infty$ y $T: X \longrightarrow X$ el operador integral definido como:
    \[ T(x)(t) = \int_{0}^{1} K(s,t) x(s) \, ds \qquad \forall x \in C^0([0,1]), \quad t \in [0,1] \]
    donde $K(t,s)$ es una función continua definida en $[0,1]^2$. Demostrar que $T$ es un operador compacto.
\end{ejercicio}

\begin{ejercicio}
    Sea $T: \ell_2 \longrightarrow \ell_2$ definido por:
    \[ T(x) = \left( x_1, \frac{x_2}{2}, \dots, \frac{x_k}{k}, \dots \right) \qquad \forall x = (x_1, \dots, x_k, \dots) \in \ell_2 \]
    Demostrar que $T$ es compacto. \\

    \noindent
    Primmero probemos que $T$ es lineal y continuo (= acotado):
    \begin{itemize}
        \item \tbf{Linealidad}: tomamos dos sucesiones $x, y \in \ell_2$ y dos escalares $\alpha, \beta \in \mathbb{C}$. Evaluando componente a componente
        $$T(\alpha x + \beta y) = \left( \alpha x_1 + \beta y_1, \frac{\alpha x_2 + \beta y_2}{2}, \dots, \frac{\alpha x_k + \beta y_k}{k}, \dots \right)$$
        Por las propiedades de las sucesiones, podemos separar los términos de $x$ y de $y$
        $$T(\alpha x+\beta y)= \alpha \left( x_1, \frac{x_2}{2}, \dots \right) + \beta \left( y_1, \frac{y_2}{2}, \dots \right) = \alpha T(x) + \beta T(y)$$
        Por tanto, $T$ es lineal.
        \item \tbf{Continuidad}: calculamos la norma al cuadrado de la imagen para cualquier $x \in \ell_2$
        $$\|T(x)\|_2^2 = \sum_{k=1}^{\infty} \left| \frac{x_k}{k} \right|^2 = \sum_{k=1}^{\infty} \frac{1}{k^2} |x_k|^2$$
        Como para todo $k \ge 1$ se cumple que $\frac{1}{k^2} \le 1$, podemos acotar la serie hacia arriba quitando los denominadores
        $$\|T(x)\|_2^2 \le \sum_{k=1}^{\infty} 1 \cdot |x_k|^2 = \|x\|_2^2$$
        Tomando raíces cuadradas, tenemos que $\|T(x)\|_2 \le 1 \cdot \|x\|_2$. Al estar acotado, el operador es continuo (y de paso sabemos que $\|T\| \le 1$).
    \end{itemize}
    Para demostrar $T$ compacto, buscaremos hacerlo como límite de operadores de rango finito (que son compactos). Para cada $n \in \mathbb{N}$, definamos el operador $T_n: \ell_2 \longrightarrow \ell_2$ por
    $$T_n(x) = \left( x_1, \frac{x_2}{2}, \dots, \frac{x_n}{n}, 0, 0, 0, \dots \right)$$
    Sabemos que son operadors de rango finito, porque para cualquier $x$, la imagen $T_n(x)$ tiene todas sus componentes nulas a partir de la posición $n+1$. Por tanto, la imagen está contenida en el subespacio generado por los primeros $n$ vectores de la base canónica ($\text{span}\{e_1, \dots, e_n\}$), que tiene dimensión exacta $n < \infty$. Está claro que es continuo (acotado), por la misma razón que $T$. \\ \\
    \noindent
    Estudiemos qué le hace el operador resta $(T - T_n)$ a un vector $x$
    $$(T - T_n)(x) = T(x) - T_n(x) = \left( 0, 0, \dots, 0, \frac{x_{n+1}}{n+1}, \frac{x_{n+2}}{n+2}, \dots \right)$$
    Las primeras $n$ componentes se cancelan por completo, y solo sobrevive la ``cola'' infinita de la sucesión. Calculamos la norma al cuadrado de esta diferencia
    $$\|(T - T_n)(x)\|_2^2 = \sum_{k=n+1}^{\infty} \frac{1}{k^2} |x_k|^2 \overset{(*)}{\le}\frac{1}{(n+1)^2} \sum_{k=n+1}^{\infty} |x_k|^2=\frac{1}{(n+1)^2} \|x\|_2^2$$
    donde en $(*)$ hemos usado que el sumatorio empieza en $k = n+1$, por lo que el denominador más pequeño de toda la serie es $(n+1)^2$. Por lo tanto, para cualquier $k$ en ese sumatorio, se cumple que 
    $$\frac{1}{k^2} \le \frac{1}{(n+1)^2}$$
    Tomando raíces cuadradas a ambos lados, tenemos que
    $$\|(T - T_n)(x)\|_2 \le \frac{1}{n+1} \|x\|_2$$
    Por la definición de norma de un operador como el supremo, deducimos que la norma del operador resta cumple
    $$\|T - T_n\| \le \frac{1}{n+1}\implies \lim_{n \to \infty} \|T - T_n\| \le \lim_{n \to \infty} \frac{1}{n+1} = 0$$
    y como la distancia en norma entre $T$ y los operadores de rango finito $T_n$ tiende a cero, concluimos que $T$ es el límite uniforme de operadores de rango finito, es decir, $T$ es compacto.
    
\end{ejercicio}

\begin{ejercicio}
    Sean $X, Y, Z$ espacios de Banach y $T, S$ dos aplicaciones lineales con $T \in BL(X,Y)$ y $S \in BL(Y,Z)$. Demostrar que si una de las dos aplicaciones entre $T$ y $S$ es compacta, entonces la composición $S \circ T$ es compacta $\big( S \circ T \in \mathcal{K}(X,Z) \big)$. \\

    \noindent
    Sea $\{x_n\} \subset X$ una sucesión acotada cualquiera en $X$. Nuestro objetivo final es extraer una subsucesión tal que $\{S(T(x_{n_{k}}))\}$ converja fuertemente en $Z$. Para ello, tenemos dos casos posibles, dependiendo de cual de las dos aplicaciones sea compacta:
    \begin{itemize}
        \item \tbf{$T$ compacto}: como la sucesión de partida $\{x_n\} \subset X$ es acotada y $T \in \mathcal{K}(X,Y)$ es compacto, la definición nos asegura que existe una subsucesión $\{x_{n_k}\}$ tal que sus imágenes por $T$ convergen fuertemente a algún elemento $y_0 \in Y$
        $$T(x_{n_k}) \longrightarrow y_0 \quad \text{en norma de } Y \text{ (cuando } k \to \infty)$$
        Como $S$ es continua, sabemos que mantiene respecta los límites fuertes, por lo que tenemos
        $$S(T(x_{n_k}))=(S\circ T)(x_{n_k}) \longrightarrow S(y_0) \quad \text{en norma de } Z$$
        y por tanto $S \circ T$ es compacto.

        \item \tbf{$S$ compacto}: como $T$ es un operador acotado (= continuo), por definición transforma conjuntos acotados en conjuntos acotados. Esto significa que la sucesión de sus imágenes en $Y$, que llamaremos 
        $$y_n = T(x_n)$$ 
        es una sucesión acotada en $Y$. Como $S$ es un operador compacto, y tenemos $\{y_n\}$ una sucesión acotada, tenemos entonces que podemos extraer una subsucesión $\{y_{n_k}\}$ tal que sus imágenes por $S$ convergen fuertemente a un elemento $z_0 \in Z$
        $$S(y_{n_k}) \longrightarrow z_0 \quad \text{en norma de } Z \text{ (cuando } k \to \infty)$$
        que reescribiendolo queda como
        $$S(T(x_{n_k})) \longrightarrow z_0 \implies (S \circ T)(x_{n_k}) \longrightarrow z_0 \quad \text{en norma de } Z$$
        y por tanto $S\circ T$ es compacto.
    \end{itemize}
\end{ejercicio}

\begin{ejercicio}
    Sea $H$ un espacio de Hilbert y $T \in BL(H)$. Entonces 
    \begin{equation*}
        T \text{ compacto} \iff T^* \text{ compacto}
    \end{equation*}
    \begin{description}
        \item[$\Longleftarrow)$] Supongamos demostrada la implicación $\Longrightarrow)$ entonces tenemos que si $T^*$ es compacto, entonces 
        $$(T^*)^* \overset{(*)}{=} T$$
        también es compacto, donde en $(*)$ hemos usado que el operador de adjunción es involutivo, es decir, que aplicarlo dos veces seguidas nos devuelve al operador original. Por lo tanto, $T$ es compacto.

        \item[$\Longrightarrow)$] Para probar que $T^*$ es compacto tomamos una sucesión acotada cualquiera $\{x_n\} \subset H$ (es decir, $\|x_n\| \le M$ para todo $n$), y nuestro objetivo es demostrar que la sucesión de imágenes por el adjunto, $\{T^*(x_n)\}$, contiene una subsucesión que converge fuertemente (en norma) en $H$. \\

        \noindent
        Es claro que como $T^*\in BL(\mathcal{H})$ es un operador acotado, entonces la sucesión de imágenes $\{T^*(x_n)\}$ también es acotada, porque para todo $n$ se cumple que
        $$\|T^*(x_n)\| \le \|T^*\| \cdot \|x_n\| \le \|T^*\| \cdot M$$
        Por lo tanto, $\{T^*(x_n)\}$ es una sucesión acotada en $H$, por lo que sabiendo que $T$ es compacto, tendríamos que para la imagen por dicho operador $\{T(T^*(x_n))\}$ existe una subsucesión, que denotaremos mediante los índices $n_k$, tal que converge fuertemente a algún elemento $z \in H$
        $$T(T^*(x_{n_k})) \longrightarrow z \quad \text{(en norma de } H \text{ cuando } k \to \infty)$$
        Como $\mathcal{H}$ es un espacio completo (de Hilbert), tenemos que toda sucesión convergente es una sucesión de Cauchy, por lo que esa subsucesión de imágenes dobles cumple que
        \begin{equation}\label{eq:limite}
            \|T(T^*(x_{n_k} - x_{n_m}))\|\overset{(*)}{=}\|T(T^*(x_{n_k})) - T(T^*(x_{n_m}))\|  \longrightarrow 0 \quad \text{cuando } k, m \to \infty
        \end{equation}
        donde en $(*)$ hemos usado la linealidad de los operadores $T, T^*$. Ahora nos centraremos en ver si la subsucesión para nuestro operador $T^*$ es de Cauchy, es decir, veámos la distancia entre dos cuales quiera términos
        \begin{align*}
            \|T^*(x_{n_k}) - T^*(x_{n_m})\|^2 &= \|T^*(x_{n_k} - x_{n_m})\|^2=( T^*(x_{n_k} - x_{n_m}), T^*(x_{n_k} - x_{n_m}) ) = \\
            \overset{(*)}{=}& ( x_{n_k} - x_{n_m}, T(T^*(x_{n_k} - x_{n_m})) ) \overset{(**)}{\le} \|x_{n_k} - x_{n_m}\| \cdot \|T(T^*(x_{n_k} - x_{n_m}))\|
        \end{align*}
        donde en $(*)$ hemos usado la propieda del operador adjunto, y en $(**)$ la desigualdad de Cauchy-Schwarz. Viendo la desigualdad triangular del primero término del producto final tenemos
        \begin{equation*}
            \|x_{n_k} - x_{n_m}\| \le \|x_{n_k}\| + \|x_{n_m}\| \le M + M = 2M
        \end{equation*}
        y entonces tenemos
        $$\|T^*(x_{n_k} - x_{n_m})\|^2 \le 2M \cdot \|T(T^*(x_{n_k} - x_{n_m}))\|$$
        Por la Ecuación (\ref{eq:limite}), el bloque de la derecha tiende a cero cuando $k, m \to \infty$, por lo que
        $$\|T^*(x_{n_k}) - T^*(x_{n_m})\|^2=\|T^*(x_{n_k} - x_{n_m})\|^2 \longrightarrow 0 \quad \text{cuando } k, m \to \infty$$
        que demuestra que la subsucesión $\{T^*(x_{n_k})\}$ es de Cauchy, y por ser $H$ un espacio completo, esa subsucesión converge fuertemente a algún elemento de $H$. Por lo tanto, $T^*$ es compacto.
    \end{description}
\end{ejercicio}

\begin{ejercicio}
    Sean $X$ e $Y$ espacios de Banach y $T \in BL(X,Y)$. Demostrar que si $T$ es compacto, entonces:
    \[ x_n \rightharpoonup x_0 \quad \text{en } X \implies T(x_n) \longrightarrow T(x_0) \quad \text{en } Y \] 

    \noindent
    Como la sucesión $\{x_n\}$ converge débilmente a $x_0$, por una consecuencia directa del Principio de Acotación Uniforme (PAU), sabemos que la sucesión es acotada en norma
    $$ \exists M > 0 \text{ tal que } \|x_n\| \le M \quad \forall n \in \mathbb{N} $$
    Supongamos, por el contrario, que la sucesión de imágenes $\{T(x_n)\}$ no converge fuertemente a $T(x_0)$. Por la definición de límite, si una sucesión no converge a un punto, significa que existe un radio $\varepsilon > 0$ y una subsucesión $\{x_{n_k}\}$ tal que sus imágenes se quedan siempre fuera de la bola de radio $\varepsilon$ centrada en $T(x_0)$
    $$ \|T(x_{n_k}) - T(x_0)\| \ge \varepsilon \quad \forall k \in \mathbb{N} $$
    El objetivo será aplicando la unicidad del límite de la convergencia débil y su relación con la convergencia fuerte, que demostremos que dicha sucesión si converge a $T(x_0)$, llegando a una contradicción, veámoslo. \\

    \noindent
    Ahora miremos la subsucesión $\{x_{n_k}\}$. Como es una subsucesión de una sucesión acotada, $\{x_{n_k}\}$ también es acotada. Como el operador $T$ es compacto, entonces de $\{x_{n_k}\}$ podemos extraer una sub-subsucesión tal que sus imágenes convergen fuertemente a algún elemento $y \in Y$
    $$ T(x_{n_{k_j}}) \longrightarrow y \quad \text{(cuando } j \to \infty) $$
    Como la convergencia fuerte implica la débil, tenemos que
    \begin{equation}\label{eq:2}
        T(x_{n_{k_j}}) \rightharpoonup y \quad \text{débilmente en } Y 
    \end{equation}
    Por otro lado, recordemos que la sucesión original cumplía que $x_n \rightharpoonup x_0$ débilmente, por lo que la subsucesión $\{x_{n_{k_j}}\}$ también conserva el mismo límite débil
    $$ x_{n_{k_j}} \rightharpoonup x_0 \quad \text{débilmente en } X $$
    y aplicando que $T$ aparte de ser compacto, es un operador lineal y continuo, entonces
    \begin{equation}\label{eq:1}
        T(x_{n_{k_j}}) \rightharpoonup T(x_0) \quad \text{débilmente en } Y
    \end{equation}
    Ahora juntando Ecuación (\ref{eq:1}) y (\ref{eq:2}), tenemos que como el límite de la convergencia débil es único, entonces
    \begin{equation*}
        T(x_0)=y\implies T(x_{n_{k_j}}) \longrightarrow y=T(x_0)
    \end{equation*}
    pero nosotros comenzamos diciendo que todos los términos de la subsucesión $\{T(x_{n_k})\}$ tenían que estar a una distancia mayor o igual que $\varepsilon$ de $T(x_0)$. Es imposible que una sub-subsucesión de ella converja a $T(x_0)$ si todos sus términos están ``lejos'' a una distancia fija $\varepsilon$, por lo que llegamos a una contradicción. \\

    \noindent
    Como asumir que la sucesión no convergía fuertemente nos ha llevado a una contradicción matemática insalvable, nuestro supuesto inicial debe ser falso. Por lo tanto, obligatoriamente
    $$ T(x_n) \longrightarrow T(x_0) \quad \text{fuertemente en norma en } Y $$
    como se quería demostrar.
\end{ejercicio}



\begin{ejercicio}
    Sean $X$ e $Y$ espacios de Banach y $T \in BL(X,Y)$. Demostrar que si $X$ es reflexivo y para cada sucesión $\{x_n\} \subset X$ se cumple la implicación:
    \[ x_n \rightharpoonup x_0 \implies \{T(x_n)\} \text{ es relativamente compacta} \]
    entonces $T$ es un operador compacto. \\

    \noindent
    Aplicando el Teorema de Banach-Alaoglu al espacio de Banach $X^*$, obtenemos que la bola cerrada $B_{X^{**}}$ es débilmente$-*$ compacta en $X^{**}$. Como $X$ es reflexivo, la inyección canónica $J: X \to X^{**}$ es sobreyectiva, por lo que identifica exactamente la bola original con la del bidual ($J(B_X) = B_{X^{**}}$), entonces como $X$ es reflexio tenemos que
    \begin{equation*}
        \text{convergencia débil y convergencia débil-* coinciden en } X^* \implies B_X \text{ es débilmente compacta en } X
    \end{equation*}
    Llegados a este punto queremos probar que $T$ es compacto, para ello tomaremos una sucesión acotada cualquier $\{x_n\} \subset X$, y veremos si de $\{T(x_n)\}$ podemos extraer una subsucesión convergente. \\

    \noindent
    Como la sucesión está acotada, eso signficia que existe un radio $R > 0$ tal que $\|x_n\| \le R$ para todo $n \in \mathbb{N}$. Por lo tanto, toda la sucesión está contenida en una bola cerrada acotada (que formalmente es un escalado de la bola unidad, $R \cdot B_X$). Como acabamos de demostrar minuciosamente usando la reflexividad de $X$ y el Teorema de Banach-Alaoglu, la bola cerrada en $X$ es débilmente compacta, lo que significa que $\{x_n\}$ contiene una subsucesión débilmente convergente en la bola cerrada. \\

    \noindent
    Por lo tanto, de nuestra sucesión acotada $\{x_n\}$ podemos extraer una subsucesión $\{x_{n_k}\}$ que converge débilmente a un elemento $x_0 \in R\cdot B_X\subseteq X$
    $$ x_{n_k} \rightharpoonup x_0 \overset{(*)}{\implies} \{T(x_{n_k})\} \text{ es relativamente compacta en } Y$$
    donde en $(*)$ hemos usado la hipótesis que nos da el enunciado. Por definición, que el conjunto de términos de la sucesión $\{T(x_{n_k})\}$ sea relativamente compacto en un espacio métrico completo (como el espacio de Banach $Y$) significa que su clausura es compacta. \\ 

    \noindent
    La consecuencia directa de esto en términos de sucesiones es que de esa misma sucesión de imágenes se puede volver a extraer una subsucesión que converge fuertemente (en norma) hacia un punto de $Y$. Denotemos a esta nueva sub-subsucesión mediante un doble subíndice como $\{x_{n_{k_j}}\}$. Entonces, existe un elemento $y \in Y$ tal que
    $$ T(x_{n_{k_j}}) \longrightarrow y \quad \text{(en norma de } Y\text{ cuando } j \to \infty) $$
    Hemos partido de una sucesión acotada cualquiera $\{x_n\} \subset X$ y, gracias a la reflexividad y a la hipótesis de compacidad relativa, hemos sido capaces de encontrar una subsucesión $\{x_{n_{k_j}}\}$ tal que sus imágenes por el operador, $\{T(x_{n_{k_j}})\}$, convergen fuertemente en norma.
\end{ejercicio}

